\documentclass[12pt]{article}

\usepackage[paperwidth=612pt,paperheight=792pt,top=72pt,bottom=72pt,left=30mm,right=20mm]{geometry}
\usepackage{times}
\usepackage{setspace}
\usepackage{fancyhdr}
\usepackage{ragged2e}
\usepackage{enumitem}
\usepackage{titlesec}
\setstretch{1.5}

\pagestyle{fancy}
\fancyhf{}
\fancyhead[R]{AyurSpace}
\fancyfoot[L]{Pillai HOC College of Engineering \& Technology (Autonomous), B. E. Computer Engineering}
\fancyfoot[R]{\thepage}
\renewcommand{\headrulewidth}{1pt}
\renewcommand{\footrulewidth}{1pt}

\setcounter{page}{1}

\begin{document}

\vspace*{\fill}
\begin{center}
\textbf{\Huge{Chapter 1}}\\
\vspace{5mm}
\textbf{\Huge{Introduction}}
\end{center}
\vspace*{\fill}
\newpage

\section*{1.1 \quad Background and Motivation}
\addcontentsline{toc}{section}{1.1 Background and Motivation}
\justify

The integration of mobile computing, artificial intelligence (AI), and conventional medicine
signifies a pivotal frontier in contemporary digital health. The World Health Organization (WHO)
says that 88\% of all countries use traditional and complementary medicine (T\&CM), with more than
170 member states reporting the use of herbal medicines, acupuncture, yoga, indigenous therapies,
and other types of traditional medicine. Ayurveda, or ``Science of Life,'' has been practiced in
India for more than 3,000 years. It takes a holistic view of health and divides human physiology
into three bio-energetic forces, or \textit{Doshas}: \textit{Vata} (Kinetic Energy),
\textit{Pitta} (Transformative Energy), and \textit{Kapha} (Cohesive Energy).

\textit{Dravyaguna Vigyan}, the study of the properties of medicinal substances, is the most
important part of Ayurvedic pharmacology. Ayurveda classifies plants based on their \textit{Rasa}
(Taste), \textit{Guna} (Quality), \textit{Virya} (Potency), \textit{Vipaka} (Post-digestive
effect), and \textit{Prabhava} (Specific action), while Western pharmacology looks for active
chemical compounds like alkaloids and glycosides. \textit{Ocimum sanctum} (Tulsi) is not only an
antibacterial agent; it is categorized as possessing a \textit{Katu} (Pungent) \textit{Rasa} and
an \textit{Ushna} (Heating) \textit{Virya}, rendering it suitable for balancing \textit{Kapha} and
\textit{Vata} disorders while potentially exacerbating high \textit{Pitta} conditions.

Even though Ayurveda is very deep, it is going through a crisis in the 21st century. The ability
to identify medicinal plants, which was once passed down through oral traditions
(\textit{Gurukula}), is fading. Urbanization has separated most people from their natural
surroundings. Because there are not enough qualified \textit{Vaidyas} (Ayurvedic doctors), millions
of people cannot get personalized diagnosis. Because of this, we need a ``Digital Vaidya'' right
away---a system that would use mobile technology to make this expert knowledge available to everyone.

\newpage

\section*{1.2 \quad Problem Statement}
\addcontentsline{toc}{section}{1.2 Problem Statement}

The digitization of Ayurveda introduces specific computational challenges that current solutions
do not adequately resolve:

\begin{enumerate}[label={\textbf{\arabic*.}}, leftmargin=*]
    \item \textbf{Taxonomic Ambiguity}: Numerous distinct species possess identical nomenclature.
    For example, ``Brahmi'' can mean either \textit{Bacopa monnieri} or
    \textit{Centella asiatica}, depending on where you are. A search engine that only uses text
    is not enough and is dangerous.

    \item \textbf{Visual Similarity}: Many plants that are good for you look like plants that are
    bad for you. For instance, the Solanaceae family has both poisonous plants and edible
    vegetables. To tell these apart, you need computer vision that works well.

    \item \textbf{Contextual Disconnection}: Linnaean taxonomy is available in current plant
    identification apps (like PlantNet), but Ayurvedic context is not. It does not help a layperson
    to know that a plant is \textit{Tinospora cordifolia} if they do not also know its
    \textit{Guduchi} properties and dosage.

    \item \textbf{Generative Hallucinations}: Large Language Models (LLMs) like GPT-4 can give you
    Ayurvedic advice, but they are prone to ``hallucinations,'' which means they make up citations
    or properties. Using an LLM for direct visual identification is presently unreliable for
    essential medical safety.
\end{enumerate}

To address these challenges, we present \textbf{AyurSpace}, a comprehensive mobile framework. Our
contributions are:

\begin{itemize}
    \item \textbf{Hybrid Neuro-Symbolic Architecture}: We put forth a new pipeline that separates
    the ``Identification'' task (using specialized CNNs) from the ``Reasoning'' task (using Semantic
    LLMs). This lessens hallucination risks while getting the most out of the context.

    \item \textbf{Digitized Dravyaguna Ontology}: We make the properties of medicinal plants into
    a searchable, object-oriented schema that makes it possible to filter with algorithms and
    perform safety checks, like ``Find plants that are safe for pregnancy.''

    \item \textbf{Quantitative Tridosha Assessment}: We translate the subjective
    \textit{Prakriti Pariksha} (Constitution Examination) into a discrete-math scoring algorithm
    that can be used over and over again, giving users the ability to check their own bio-energy
    level.

    \item \textbf{Open Source Mobility}: We provide a reference implementation in Flutter, making
    sure the solution works on all platforms, is fast, and can be used on cheap devices common in
    developing countries.
\end{itemize}

\newpage

\section*{1.3 \quad Organisation of Report}
\addcontentsline{toc}{section}{1.3 Organisation of Report}
\justify

This report is composed of the following sections:

\begin{description}
    \item[\textbf{Chapter 1:}] Introduction\\
    This chapter introduces the AyurSpace project, highlighting the background and motivation of
    modernizing Ayurveda and providing a digital solution, and gives an overview of the report's
    structure.

    \item[\textbf{Chapter 2:}] Literature Review\\
    This chapter surveys existing systems, methods, and literature in botanical computer vision,
    LLMs in healthcare, and digital approaches to Ayurveda.

    \item[\textbf{Chapter 3:}] Requirement Gathering\\
    This chapter lists the software and hardware requirements required for the project.

    \item[\textbf{Chapter 4:}] Plan of Project\\
    A project plan outlines the approach, tasks, timelines, resources, and deliverables for
    completing the AyurSpace application.

    \item[\textbf{Chapter 5:}] Project Analysis\\
    Project analysis involves the systematic examination and evaluation through use case diagrams,
    class diagrams, and sequence diagrams.

    \item[\textbf{Chapter 6:}] Project Design\\
    This chapter details the data flow diagrams (level 0 and 1) and flow charts that map the
    application's processing sequence.

    \item[\textbf{Chapter 7:}] Implemented System\\
    This chapter details the system architecture (Clean Architecture, API Integration) and reveals
    sample code from the implemented system.

    \item[\textbf{Chapter 8:}] Result Analysis\\
    This chapter shows the performance benchmarks, accuracy metrics, latency breakdowns, and UI
    snapshots of AyurSpace.

    \item[\textbf{Chapter 9:}] Conclusion and Future Scope\\
    This chapter summarizes the conclusion of this research work with the future scope for edge AI
    and telemedicine integration.
\end{description}

\end{document}
